\section{Introduction}
\label{sec:introduction}

The stress that drives plastic flow in a saturated porous solid is not the
stress that a drained experiment measures.  A drained test loads the skeleton
at fixed pore pressure and records the Biot effective stress
\(\bs\sigma^{\prime\prime}\), the effective stress held at fixed pressure; the
dissipation of inelastic deformation is conjugate to a different stress, the
Terzaghi effective stress \(\bs\sigma^{\prime}\), held at fixed intrinsic
density.  Foster and Xu identify the two as the fixed-pressure and
fixed-density effective stresses of a finite-deformation solid--water mixture
and connect them through the Biot coefficient,
\(\bs\sigma^{\prime}=\bs\sigma^{\prime\prime}+(1-B)p\mbf I\)
\citep{fosterxu2025}: the Biot effective stress is the stress a drained test
controls, and the Terzaghi effective stress is the stress conjugate to
inelastic deformation \citep{drumheller2000}.  The two coincide only when
\(B=1\), the limit of an
incompressible mineral; for a compressible mineral the difference
\((1-B)p\mbf I\) separates the drained stress from the plastic driving
stress.  A yield function and a flow rule for such a solid therefore cannot
be evaluated until the coefficient \(B\) is known.  We realize this coefficient
in a monolithic finite-element system, show that it is a history-dependent
tangent of the active constitutive system rather than an elastic modulus
ratio, and measure the consequences for the predicted plastic response.

The consequences are not second order.  Activating an isochoric true-plastic
mechanism together with a scalar plastic pore allocation \(a^p\) raises the
coefficient toward unity: in a drained single-point compression the corrected
coefficient exceeds the elastic value \(B_{\mathrm{el}}\) by \(0.003\) at
\(5\%\) compression and by \(0.043\) at \(30\%\), over a range in which the
elastic coefficient itself falls with the finite-strain response it
quantifies.  A coefficient held at its elastic value therefore over-predicts
the plastic increment, by about three percent at \(400\,\mathrm{MPa}\) pore
pressure and \(20\%\) compression, and the error grows with pore pressure and
with compression.  The coefficient also moves before yield: in the
finite-deformation Mandel continuation the spatially averaged coefficient
decreases from \(0.600\) to \(0.51417\), a relative change of \(14.3\%\), over
ten percent platen compression.  Each of these effects follows from the same
implicit, history-dependent construction, and none is expressible by a
constant coefficient.

The finite-deformation coefficient measures the change of mineral distension
with skeleton Jacobian at fixed pore pressure
\citep{fosterxu2025}.  In the small-strain limit it reduces to the familiar
drained-modulus ratio of poroelasticity \citep{biotwillis1957}.  The
one-solid, one-water specialization retains solid-reference kinematics, phase
material measures, and the pressure Legendre transform, so the coefficient is
a fixed-pressure constitutive derivative.  Solid mass conservation is a
spatial balance for the bulk solid partial density.  The local constitutive
system supplies the intrinsic solid state, and the solid volume fraction is
recovered algebraically from the partial and intrinsic densities.  The
fixed-pressure derivative of the local mineral distension then supplies the
Biot coefficient; it is evaluated by implicit differentiation of the
constitutive residual system while holding intrinsic phase pressure and
referential solid mass \(J\rho_s\) fixed.  The coefficient enters the stress
residual, and the derived solid volume fraction enters the complete nonlinear
water accumulation, with both paths active in the global Newton linearization.

The new content is the monolithic finite-element realization and verification
of this sequence for a finite-deformation solid--water system.  Sun, Ostien,
and Salinger show how residual evaluators templated on scalar type provide
automatic derivatives for a coupled finite-strain poromechanics Jacobian
\citep{sunostien2013}; the Multiphysics Object Oriented Simulation Environment
(MOOSE) applies the same residual-based approach through automatic
differentiation (AD) \citep{gastonetal2009moose,lindsayetal2021ad}, and
consistent algorithmic tangents provide the corresponding link between a local
material update and the global Newton method
\citep{simotaylor1985,miehe1996}.  Large-deformation porous-material
formulations likewise separate skeleton deformation, fluid transport, and
constitutive state evolution
\citep{hongetal2008,macminnetal2016,dehghanizilian2021}.  We implement the
residual-vector architecture required by a coupled nonlinear local update:
after the local state converges, an AD-valued dense solve differentiates the
active constitutive system and returns its consistent fixed-pressure Biot
tangent.  The elastic specialization has a closed form and is used only to
verify this general path.  AD retains the complete outer dependence without a
hand-coded global tangent, and the same residual-vector interface admits
inelastic internal variables and an active-set-consistent tangent without
assigning a balance law to a constitutive state.

The distinction between the two effective stresses becomes decisive once the
skeleton yields.  In poroplasticity the yield function and plastic potential
are written in the Terzaghi effective stress \(\bs\sigma^{\prime}\); the
dissipation analysis identifies the driving stresses for plasticity
\citep{drumheller2000}: in the multicomponent reduction the isochoric
true-plastic mechanism is driven by the deviatoric part of
\(\bs\sigma^{\prime}\) and the pore-allocation mechanism by its mean part
\citep{fosterxu2025}.  Reconstructing that driving stress from the Biot
effective stress through
\(\bs\sigma^{\prime}=\bs\sigma^{\prime\prime}+(1-B)p\mbf I\) requires
the coefficient \(B\), which is itself the fixed-pressure tangent of the
active local system.  The coefficient and the plastic state are therefore
solved together: \(B\) is needed to define the driving stress, and the plastic
state is needed to compute \(B\).  An implicit, history-dependent \(B\) is
therefore not incidental bookkeeping but the object that makes a
poroplasticity model with a compressible mineral well posed.  We give that
construction for a solid whose inelastic response combines an isochoric
true-plastic mechanism with a scalar plastic pore allocation, verify its
elastic limit against the closed forms and the Mandel benchmark, and report
the coefficient correction together with its feedback through the
\((1-B)p\) term at nonzero pore pressure.  The history dependence parallels
the stress-dependent coefficient measured in drained compactive sandstone
\citep{ingraham2017evolution}.

This paper has four objectives.  First, it states the single-solid,
single-water finite-deformation equations used by the implementation, the
spatial solid and water mass balances, the elastic mineral and water equations
of state, and the fixed-pressure Biot transform.  Second, it shows how the
local constitutive outputs remain active in the finite-element Jacobian
through the implicit-tangent construction of
Sec.~\ref{sec:implicit-ad-implementation}.  Third, it verifies the implemented
residual path with the Mandel problem by comparing water-pressure and
displacement profiles with the analytical plane-strain solution of Cheng and
Detournay \citep{chengdetournay1988}, and it reports the finite-deformation
continuation in which the coefficient shifts from \(0.600\) to \(0.51417\).
Fourth, it develops the inelastic construction: an isochoric true-plastic
mechanism together with a scalar plastic pore allocation, whose driving stress
is the Terzaghi effective stress reconstructed from the Biot effective stress
through the coefficient returned by the active implicit solve.  The elastic closure
verifies the inactive-mechanism limit, and the material-level demonstration
reports the coefficient correction and its feedback through the \((1-B)p\)
term.  That demonstration does not reproduce a particular rock's drained
behavior, which would require a separate constitutive ingredient outside the
present scope.
